\begin{lemma}
Let $H\colon\mathbb{C}\to\mathbb{C}$, let $\Omega\subseteq\mathbb{C}$ be open
with $H$ of class $C^\infty$ on $\Omega$ and $\xi\neq0$ for every
$\xi\in\Omega$, let $g\in\mathbb{R}$, and let $C_m$, $m=0,1,2,\dots$, be reals
such that
\[
  \sum_{|w|=m}\bigl|D^{w}H(\xi)\bigr|\le C_m|\xi|^{g-m}
  \qquad\text{for every }m\ge0\text{ and every }\xi\in\Omega,
\]
the sum running over the $2^m$ words $w$ of length $m$ in $\{0,1\}$ and $D^{w}$
being the iterated partial derivative of \noderef{PartialWord}.  Let
$s,k\ge0$ be integers and let $P\in\mathbb{P}(s,k)$ relative to $H$ in the
sense of \noderef{PClass}.  Then there is $M\ge0$ with
\[
  \bigl|P(\xi)\bigr|\le M\,|\xi|^{kg-s}\qquad\text{for every }\xi\in\Omega .
\]

This is the growth estimate the paper records at paper.tex:488 for a monomial
$Q=H^{s_1}\cdots H^{s_k}$ of type $(s,k)$, namely
$|Q(\xi)|\le M|\xi|^{g-s_1}\cdots|\xi|^{g-s_k}=M|\xi|^{kg-s}$, extended to all
of $\mathbb{P}(s,k)$ by summing over the finite linear combination that
membership in $\mathbb{P}(s,k)$ carries.  Applied with $g=\gamma_n$, the
hypothesis is estimate (43) of Corollary 1 and the conclusion at
$P_\nu\in\mathbb{P}(4l-3,3l-2)$ is the paper's
$|P_\nu(\xi)|\le M|\xi|^{(3l-2)\gamma_n-4l+3}$ of paper.tex:492.
\end{lemma}
\begin{proof}
If $\Omega=\emptyset$ the assertion holds with $M=0$, the conclusion being
vacuous.  So assume $\Omega\neq\emptyset$ and fix $\xi_0\in\Omega$.

\emph{Step 1: $C_m\ge0$ for every $m$.}  The left-hand side of the hypothesis
at $\xi_0$ is a sum of moduli, hence $\ge0$, so
$0\le C_m|\xi_0|^{g-m}$.  By assumption $\xi_0\neq0$, so $|\xi_0|>0$ and
$|\xi_0|^{g-m}>0$; dividing, $C_m\ge0$.

\emph{Step 2: every Wirtinger derivative of order $q$ is controlled by the
real ones.}  We claim: if $G$ is $C^\infty$ on $\Omega$ and $T_1,\dots,T_q$ is
any list with each $T_a\in\{\partial/\partial\xi,\partial/\partial\bar\xi\}$,
then on $\Omega$
\begin{equation}\tag{1}
  T_1T_2\cdots T_q\,G=2^{-q}\sum_{|w|=q}\varepsilon_w\,D^{w}G,
\end{equation}
for some complex numbers $\varepsilon_w$ with $|\varepsilon_w|\le1$, the sum
being over the $2^q$ words of length $q$.

We induct on $q$.  For $q=0$ the composite is $G=D^{[\,]}G$, so (1) holds with
the single coefficient $1$.  For the step, let $T_1\cdots T_qG=T_1\bigl(T_2\cdots T_qG\bigr)$
and write $G'=T_2\cdots T_qG$.  By \noderef{DZ} and \noderef{DZbar},
\[
  \frac{\partial G'}{\partial\xi}
    =\tfrac12\bigl(D^{[0]}G'-i\,D^{[1]}G'\bigr),\qquad
  \frac{\partial G'}{\partial\bar\xi}
    =\tfrac12\bigl(D^{[0]}G'+i\,D^{[1]}G'\bigr),
\]
the two real partials being the values of the Fréchet derivative on $1$ and
$i$, i.e.\ $D^{[0]}$ and $D^{[1]}$ in the sense of \noderef{PartialWord}.  By
the inductive hypothesis $G'=2^{-(q-1)}\sum_{|w'|=q-1}\varepsilon_{w'}D^{w'}G$
as functions on the open set $\Omega$, where each $D^{w'}G$ is $C^\infty$ on
$\Omega$ because $G$ is; so $\partial_t$ may be applied term by term, and by
the recursion $D^{t::w'}=\partial_t D^{w'}$ of \noderef{PartialWord},
\[
  D^{[t]}G'=2^{-(q-1)}\sum_{|w'|=q-1}\varepsilon_{w'}D^{t::w'}G
  \qquad(t=0,1).
\]
Substituting into the two displays above expresses $T_1G'$ as
$2^{-q}\sum_{|w|=q}\varepsilon_wD^{w}G$ with each $\varepsilon_w$ equal to
$\pm\varepsilon_{w'}$ or $\pm i\varepsilon_{w'}$ for the corresponding $w'$,
hence of modulus at most $1$.  This proves (1).

Now let $d$ be a derivative symbol in the sense of \noderef{DerivSymbol}, with
orders $i,j$ and total order $q=i+j$; note $q\ge1$, since
\noderef{DerivSymbol} requires $0<i+j$.  By \noderef{WirtingerIter} the
iterated Wirtinger derivative $\partial^{q}H/\partial\xi^{i}\partial\bar\xi^{j}$
is the composite of $j$ copies of $\partial/\partial\bar\xi$ followed by $i$
copies of $\partial/\partial\xi$ applied to $H$, which is a composite as in
(1) with $G=H$; and by \noderef{SymbolValue} the value of $d$ on $H$ is that
iterated derivative or its complex conjugate, which has the same modulus.
Hence, for $\xi\in\Omega$, by (1) and $|\varepsilon_w|\le1$,
\[
  \bigl|H^{d}(\xi)\bigr|
  \le2^{-q}\sum_{|w|=q}\bigl|D^{w}H(\xi)\bigr|
  \le\sum_{|w|=q}\bigl|D^{w}H(\xi)\bigr|
  \le C_q|\xi|^{g-q},
\]
writing $H^{d}$ for the value of $d$ on $H$, using $2^{-q}\le1$ and all terms
nonnegative, and applying the hypothesis with $m=q$.  So
\begin{equation}\tag{2}
  \bigl|H^{d}(\xi)\bigr|\le C_{q_d}\,|\xi|^{g-q_d}\qquad(\xi\in\Omega),
\end{equation}
where $q_d\ge1$ is the order of $d$.

\emph{Step 3: the monomial bound.}  Let $\mathcal{M}=[d_1,\dots,d_k]$ be a list
of derivative symbols with $\sum_{a=1}^{k}q_{d_a}=s$, i.e.\ a monomial of type
$(s,k)$, and let $\xi\in\Omega$.  By \noderef{MonomialValue} its value is the
product $\prod_{a=1}^{k}H^{d_a}(\xi)$, so by (2), all factors and all bounds
being nonnegative by Step 1,
\[
  \bigl|Q_{\mathcal{M}}(\xi)\bigr|
  =\prod_{a=1}^{k}\bigl|H^{d_a}(\xi)\bigr|
  \le\prod_{a=1}^{k}C_{q_{d_a}}|\xi|^{g-q_{d_a}}
  =\Bigl(\prod_{a=1}^{k}C_{q_{d_a}}\Bigr)\,|\xi|^{\sum_{a}(g-q_{d_a})}
  =\Bigl(\prod_{a=1}^{k}C_{q_{d_a}}\Bigr)|\xi|^{kg-s},
\]
where the powers of the positive real $|\xi|$ were combined by
$|\xi|^{\alpha}|\xi|^{\beta}=|\xi|^{\alpha+\beta}$, legitimate because
$|\xi|>0$, and $\sum_{a}(g-q_{d_a})=kg-\sum_aq_{d_a}=kg-s$.  This is exactly
the paper's display at paper.tex:488.

\emph{Step 4: the linear combination.}  By \noderef{PClass} there is a finite
list $L$ of pairs $(c,\mathcal{M})$, with every $\mathcal{M}$ occurring in $L$
of type $(s,k)$ --- that is, of length $k$ and with orders summing to $s$ ---
such that
\[
  P(\xi)=\sum_{(c,\mathcal{M})\in L}c\,Q_{\mathcal{M}}(\xi)
  \qquad\text{for every }\xi .
\]
Hence, by the triangle inequality and Step 3, for $\xi\in\Omega$,
\[
  \bigl|P(\xi)\bigr|
  \le\sum_{(c,\mathcal{M})\in L}|c|\,\bigl|Q_{\mathcal{M}}(\xi)\bigr|
  \le\Bigl(\sum_{(c,\mathcal{M})\in L}|c|\prod_{a=1}^{k}C_{q_{d_a(\mathcal{M})}}\Bigr)
    |\xi|^{kg-s}
  =M\,|\xi|^{kg-s},
\]
where $M$ denotes the bracketed sum.  It is a finite sum of products of
nonnegative numbers by Step 1, hence $M\ge0$.  Note that the exponent is the
same for every term because \noderef{PClass} requires every monomial in the
presentation to have the \emph{same} type $(s,k)$.  This is the assertion.
\end{proof}
